% ch11.tex - 第十一章 相关工作与对比分析
\chapter{相关工作与对比}
\label{ch:related-work}

本章将本系统与工业巡检机器人现有方案进行对比，明确技术定位与差异化特征。对比数据来自公开行业资料。

\section{现有方案概况}
\label{sec:existing-solutions}

工业巡检机器人是已有成熟商业产品的领域，应用集中于电力系统巡检。据公开资料与国网、南网招标结果，代表性厂商包括亿嘉和、申昊科技、朗驰欣创、浙江国自与鲁能智能等，产品具备红外测温、局部放电检测与表计识别等功能。这类方案的共性特征是：以电力固定设施为主要场景；部分仍依赖预设轨道或固定巡检点位；产品形态为面向特定行业的封闭商业系统，其操作系统与算法体系的自主可控程度未对外公开。

\section{技术维度对比}
\label{sec:comparison-matrix}

表~\ref{tab:competitor-comparison}从十一个维度将本系统与典型电力巡检机器人方案进行对比。

\begin{table}[htbp]
\centering
\caption{本系统与现有方案的逐项对比}
\label{tab:competitor-comparison}
\small
\begin{tabularx}{\textwidth}{p{2.2cm}XX}
\toprule
\textbf{对比维度} & \textbf{现有方案（电力巡检机器人）} & \textbf{本系统} \\
\midrule
操作系统 & Linux/Android 或闭源自研，自主可控程度不透明 & OpenHarmony 5.0，国产分布式操作系统 \\
AI 算力 & 通用 GPU 或境外芯片 & 华为昇腾 NPU，国产边缘算力 \\
算法栈 & 多基于 ROS 或闭源技术栈 & 核心算法自主实现，无 ROS 依赖 \\
移动方式 & 部分依赖预设轨道或固定点位 & 无轨自主 SLAM 导航，适用于任意室内场景 \\
覆盖能力 & 多为固定点位巡检 & BCD 全覆盖遍历，适用于任意非凸形状区域 \\
多机协同 & 单机作业或中心服务器调度 & 软总线共享黑板，去中心协同 \\
大模型分析 & 云端分析或不具备 & DeepSeek 端侧离线分析，数据不出本体 \\
数据安全 & 部分依赖云端传输 & 全程设备本体离线，支持断网运行 \\
读数可解释性 & 多为端到端黑盒识别 & 两段式几何换算，可标定、可复现 \\
自主可控 & 部分环节存在外部依赖 & 内核、算力、语言、算法、硬件全栈国产化 \\
可扩展性 & 面向特定行业的封闭产品 & 模块化架构，支持多模态扩展 \\
\bottomrule
\end{tabularx}
\end{table}

\section{差异化特征}
\label{sec:differentiation}

基于上述对比，本系统的差异化特征可归纳为四点。

第一，全栈国产化与技术自主。现有方案的底层技术栈对境外操作系统、芯片与 ROS 生态的依赖程度通常不透明，存在供应链安全风险；本系统从操作系统、AI 算力、开发语言、核心算法到硬件平台形成完整的国产化路径，在信创替代的政策背景下具有体系结构层面的差异。

第二，无轨自主全覆盖。部分现有方案依赖预设轨道或固定点位，部署成本高、灵活性有限；本系统采用无轨 SLAM 导航配合 BCD 全覆盖，无需改造现场即可在任意室内场景自主建图并生成不重不漏的覆盖路径，对含柱体与设备的非凸场地同样适用。

第三，端侧感知与分析一体化。多数现有方案的仪表识别停留在“检测并上报”，后续分析依赖云端或人工；本系统通过 NPU 分时调度，在同一块昇腾芯片上同时承载实时视觉推理与离线大模型分析，数据全程不出本体。

第四，去中心化多机协同。现有多机方案多为单机作业或中心服务器调度；本系统基于鸿蒙软总线的共享黑板与车载 agent，实现了无中心服务器、由状态变化驱动的多机协同。

需要说明的是，现有商业产品经过大规模部署验证，在防护等级、行业认证与长期稳定性方面具有明显优势。本系统作为面向技术验证的原型，其价值在于验证了“全栈国产化 + 无轨全覆盖 + 端侧大模型 + 软总线协同”路线的工程可行性，为自主可控巡检机器人的产业化提供技术参考。

\section{政策契合度}
\label{sec:policy-alignment}

本系统的技术路线与多项国家政策导向契合（表~\ref{tab:policy-alignment}）。

\begin{table}[htbp]
\centering
\caption{技术路线与政策导向的契合点}
\label{tab:policy-alignment}
\small
\begin{tabularx}{\textwidth}{p{5cm}X}
\toprule
\textbf{政策导向} & \textbf{契合点} \\
\midrule
《“十四五”机器人产业发展规划》重点发展特种机器人 & 本系统面向工业巡检、危险环境作业场景 \\
2025 年政府工作报告首次提出“具身智能” & 系统集成自主移动、环境感知、视觉智能与端侧大模型 \\
信创替代：2027 年央国企关键软硬件国产化 & 全栈国产化路线可直接服务电力、能源等信创场景 \\
昇腾 CANN 开源与 OpenHarmony 生态建设 & 基于昇腾边缘算力与 OpenHarmony 完成原生部署 \\
\bottomrule
\end{tabularx}
\end{table}

综上，本系统在全栈国产化、无轨全覆盖、端侧感知--分析一体化与去中心化协同四个方面与现有方案形成明确差异，且与国家产业政策方向一致，具备进一步工程化与产业化的价值。
